\section{Experimental Design}

The experiment is designed to check whether the proposed audit workflow can be instantiated in a controllable data-circulation setting and whether the lightweight path can be compared with full-node confirmation. The design contains two parts: an end-to-end containerized workflow experiment and a mechanism-level ablation experiment.

\subsection{Dataset and Data-Circulation Workload}

The end-to-end experiment uses the standard MNIST dataset as a reproducible workload. MNIST contains 60,000 training samples and 10,000 test samples. The training data are randomly partitioned into ten participants to emulate a multi-party data-circulation process in which each participant contributes local updates rather than directly exposing raw data. The task is run for 100 training rounds. During the workflow, model updates, audit events, update digests, and abnormal-update records are treated as auditable evidence objects.

To test abnormal-event handling, malicious updates are injected at predefined rounds. In the experiment setting, rounds 17, 37, 57, 77, and 97 are marked as abnormal rounds, and the malicious update is generated by multiplying the update direction by a negative factor. This injection creates traceable abnormal events for the audit workflow.

\subsection{Containerized Blockchain-Audit Environment}

The system is deployed with Docker Compose. The containerized environment includes ZooKeeper, blockchain-audit nodes, and a runner. ZooKeeper maintains the node registry and node status. The runner reads registered and healthy nodes from ZooKeeper, constructs audit payloads, and submits them to the available audit nodes for PBFT-style voting. Each committed audit block records the payload hash, Merkle root, previous block hash, voting result, and related audit metadata.

The node-governance experiment uses four audit nodes. Three governance scenarios are considered. In the normal scenario, all four nodes are trusted and available. In the low-trust scenario, one node remains registered but its trust score is below the threshold, so it is excluded from the voting set. In the faulty-node scenario, one node is registered and trusted but unavailable during confirmation, while the remaining nodes must still satisfy the quorum rule.

\subsection{Compared Strategies}

The workflow experiment compares four strategies. The first is ASGD without chain auditing, which provides a no-chain training baseline. The second is a two-factor no-chain strategy, which adds local anomaly screening but does not commit audit evidence on chain. The third is full PBFT chain auditing, in which audit events are confirmed by the full trusted node set. The fourth is the TrustAuditFlow audit-gate strategy, in which routine audit events use trust-aware gating and chain anchoring, while abnormal or high-risk events are escalated.

The mechanism-level ablation further decomposes the blockchain-audit path into four variants: BLoP-style full-node PBFT, full-node PBFT with Merkle batch anchoring, committee confirmation without recovery closure, and the complete TrustAuditFlow path with committee confirmation, batch anchoring, recovery-state recording, and fallback.

\subsection{Metrics and Verification Procedure}

The experiment focuses on audit-workflow properties. The main metrics are: number of on-chain records, confirmation latency, quorum satisfaction, abnormal-event recording, hash-chain continuity, Merkle-root consistency, node inclusion or exclusion under trust gating, and evidence redundancy. For each run, a verifier checks whether committed blocks form a continuous hash chain, whether the expected quorum is satisfied, whether abnormal rounds are recorded, and whether each scenario produces an auditable chain state.

The experimental procedure is as follows. First, the MNIST files are prepared and mounted into the runner container. Second, ZooKeeper and audit nodes are started through Docker Compose. Third, the runner executes the selected strategy and emits audit events during training rounds. Fourth, audit payloads are confirmed by the corresponding voting set and appended to the audit chain. The verifier then checks the generated chain files, node states, abnormal-event records, and summary metrics. This procedure provides a repeatable basis for comparing full-node auditing and lightweight audit gating.
